\section{Thảo luận}
\label{sec:discussion}

\subsection{Điều gì quyết định kết quả?}

Kết quả nổi bật nhất của báo cáo không phải là một mô hình có RMSE thấp, mà là việc
một baseline đơn giản như Naive persistence lại khó bị đánh bại. Điều này không quá
ngạc nhiên nếu nhìn vào ACF: lưu lượng hôm nay tương quan rất mạnh với lưu lượng
những ngày trước, và mức nền thay đổi chậm. Trong một chuỗi như vậy, phần lớn thông
tin dự báo ngắn hạn nằm ở giá trị gần nhất; mô hình học máy chỉ có thể cải thiện nếu
nó học được phần thay đổi nhỏ vượt trên tính dai dẳng.

LightGBM được thiết kế với mục tiêu delta để tận dụng chính tính dai dẳng này. Nếu
mô hình dự báo delta bằng 0, nó trở thành Naive; nếu dự báo delta khác 0, nó cố gắng
cải thiện. Trên Validation và Test, các delta học được có ích ở một số thời điểm
nhưng không đủ để tạo ra cải thiện tổng thể. Kết quả là LightGBM gần như trùng Naive
ở D+1--D+3 và kém hơn một chút ở D+5--D+7. Đây là một kết quả \emph{âm tính có ý
nghĩa}: nó cho thấy mô hình phức tạp không tự động tốt hơn trong dữ liệu có cấu
trúc đơn giản.

Các lần thử tham số cũng cho thấy điều tương tự. Tám cấu hình khác nhau cho mean
RMSE từ 5,57 đến 5,58; khoảng cách quá nhỏ so với độ bất định của hai năm Test. Việc
tối ưu tham số không phải là nút thắt chính. Nút thắt là thông tin: chuỗi chỉ có
một biến, không có mưa, nhiệt độ, tuyết tan hay vận hành hồ chứa. Những biến này có
thể giải thích phần thay đổi mà lag và rolling window không nắm được.

\subsection{So sánh với kỳ vọng và tài liệu}

Trong lý thuyết dự báo chuỗi thời gian, các mô hình đơn giản thường là baseline bắt
buộc, đặc biệt khi chuỗi có tính dai dẳng cao \cite{hyndman2021}. Kết quả của báo cáo
phù hợp với nguyên tắc này: trước khi kết luận một mô hình học máy tốt hơn, cần so
sánh với Naive trên cùng tập đánh giá. Việc báo cáo giữ Naive trong mọi bảng so sánh
là cần thiết, không phải hình thức.

Seasonal Naive 365 thất bại vì giả định mùa vụ năm ổn định. Phân rã STL cho thấy
thành phần mùa vụ yếu hơn nhiều so với thành phần xu hướng; ACF ở độ trễ 365 ngày
còn âm nhẹ. Điều này không có nghĩa là chuỗi không có mùa vụ, mà là mùa vụ không đủ
mạnh và ổn định để dùng làm dự báo trực tiếp. Khi mức nền thay đổi giữa các năm,
việc lặp lại cùng ngày năm trước có thể tạo sai số lớn hơn cả một mô hình “không
học” gì ngoài giá trị gần nhất.

LightGBM là một triển khai gradient boosting hiệu quả \cite{ke2017}; tuy nhiên, sức
mạnh của nó phụ thuộc vào đặc trưng và dữ liệu. Với 29 đặc trưng nội tại của một
chuỗi, mô hình có thể học các tương tác phi tuyến nhưng không thể tạo ra thông tin
mới. Feature importance cho thấy các đặc trưng lịch và lag dài được dùng nhiều; điều
này hợp lý nhưng cũng cảnh báo rằng mô hình có thể đang học các quy luật mức nền
thay vì quan hệ nhân quả. Vì vậy không nên diễn giải gain như bằng chứng nhân quả.

\subsection{Giới hạn dữ liệu và thiết kế}

Báo cáo có một số giới hạn cần được nêu rõ:

\begin{itemize}[leftmargin=1.4em]
  \item \textbf{Một chuỗi, một trạm:} không thể đánh giá khả năng tổng quát hóa sang
  trạm khác hoặc sang lưu vực khác.
  \item \textbf{Không có biến ngoại sinh:} thiếu mưa, nhiệt độ, tuyết tan, bốc hơi
  và vận hành hồ chứa; đây có thể là phần lớn nguyên nhân của các thay đổi mức nền.
  \item \textbf{Test ngắn:} chỉ hai năm 2022--2023; kết quả có thể thay đổi nếu chọn
  giai đoạn Test khác.
  \item \textbf{Không có cực trị trong Test:} 2022--2023 không có ngày vượt Q95, nên
  không thể đánh giá khả năng dự báo lũ.
  \item \textbf{Ngưỡng Q95 toàn mẫu:} ngưỡng được tính trên toàn bộ 2010--2023,
  không phải ngưỡng trượt; phù hợp mô tả nhưng không phù hợp vận hành.
  \item \textbf{Không có khoảng bất định:} báo cáo chỉ đưa ra dự báo điểm, không có
  khoảng tin cậy hoặc phân phối dự báo.
  \item \textbf{Không có kiểm định ADF/KPSS:} kết luận về tính dừng dựa trên ACF và
  STL, không dựa trên kiểm định chính thức.
  \item \textbf{Không có GRU/LSTM:} môi trường không cài PyTorch; báo cáo không đưa
  ra kết quả cho các mô hình này.
\end{itemize}

Những giới hạn này không làm kết quả hiện tại mất giá trị, nhưng chúng xác định
phạm vi diễn giải. Kết quả đúng nhất khi được đọc là: “trong dữ liệu 2010--2023,
với các đặc trưng nội tại và split đã chọn, LightGBM không vượt Naive”; không phải
“LightGBM không thể dự báo chuỗi này” hay “không mô hình nào có thể tốt hơn”.

\subsection{Hướng cải thiện}

Có bốn hướng cải thiện tự nhiên. Thứ nhất, bổ sung biến ngoại sinh như mưa, nhiệt
độ, tuyết tan và thông tin vận hành hồ chứa; đây là hướng có khả năng tạo cải thiện
lớn nhất. Thứ hai, dùng đánh giá xác suất (quantile regression, conformal prediction)
để trả lời câu hỏi “dự báo có thể sai bao nhiêu”, thay vì chỉ một con số điểm. Thứ
ba, thử nghiệm ngưỡng cực trị trượt và đánh giá riêng trên các đợt cực trị. Thứ tư,
mở rộng sang nhiều trạm hoặc nhiều lưu vực để kiểm tra khả năng tổng quát hóa.

Với mô hình, có thể thử các biến thể như LightGBM với monotonic constraints, mô hình
tuyến tính có lag, hoặc mô hình trạng thái (state-space). Tuy nhiên, mọi biến thể cần
được so sánh với Naive trên cùng split và cùng horizon. Nếu không có cải thiện rõ
ràng, giữ baseline đơn giản là lựa chọn hợp lý hơn về chi phí vận hành và khả năng
giải thích.

\subsection{Ý nghĩa thực hành}

Trong bối cảnh ra quyết định, kết quả này nhắc rằng một mô hình phức tạp cần được
chứng minh bằng thực nghiệm, không nên giả định trước. Với chuỗi lưu lượng có tính
dai dẳng cao, dự báo ngắn hạn có thể dùng persistence làm mốc; mô hình học máy chỉ
nên được thêm vào nếu có bằng chứng cải thiện ngoài mẫu. Đồng thời, cần theo dõi
mức nền và các giai đoạn chuyển chế độ, vì đó là nơi các mô hình dựa trên mùa vụ dễ
sai nhất.

Cuối cùng, việc công bố mã nguồn, dữ liệu và website giúp người đọc kiểm tra lại
các kết luận. Đây là một phần của chất lượng báo cáo: không chỉ trình bày kết quả
mà còn cho thấy cách tạo ra kết quả đó.
